% ============================================================
%  EXAMEN TEMA 1: RAZONAMIENTO LÓGICO Y LENGUAJE MATEMÁTICO
%  Curso Introductorio — 3er Año
%  Duración: 90 minutos
%  Temas: Lógica proposicional, Método de Pólya y
%         traducción al lenguaje algebraico
%  Autor: Ing. Gabriel Astudillo
%  Nota: enunciado limpio (sin pistas ni desarrollo). Las
%  soluciones están en la "Clave de Respuestas" al final.
% ============================================================

\documentclass[12pt, letterpaper]{article}

% ── Paquetes ─────────────────────────────────────────────────

\usepackage{mathwizards-guia}
\mwguia{Examen — Tema 1: Razonamiento Lógico y Lenguaje Matemático}{Ing. Gabriel Astudillo $\cdot$ Curso 3er Año}

\begin{document}
% ============================================================

% ── Portada / Título ─────────────────────────────────────────
\mwportada{Examen del Tema 1}{Razonamiento Lógico y Lenguaje Matemático}{Lógica Proposicional, Método de Pólya y Traducción Algebraica}

\vspace{4pt}
\begin{cuadroinfo}
  \small
  \textbf{Nombre:} \rule{0.55\linewidth}{0.4pt} \quad \textbf{Fecha:} \rule{0.15\linewidth}{0.4pt} \\[6pt]
  \textbf{Tiempo disponible:} 90 minutos \quad$\bullet$\quad \textbf{Puntaje total:} 100 puntos \\[4pt]
  \textbf{Instrucciones:}
  \begin{enumerate}[leftmargin=*]
    \item Lee cada ejercicio con calma antes de resolverlo.
    \item Justifica \emph{todos} tus razonamientos; en los acertijos, explica el paso a paso.
    \item En los problemas de Pólya, escribe explícitamente las cuatro fases.
    \item Define las proposiciones atómicas que uses antes de traducir.
    \item No se permite el uso de calculadora.
    \item Escribe con letra clara y revisa tu examen antes de entregarlo.
  \end{enumerate}
\end{cuadroinfo}

\vspace{8pt}

% ============================================================
\section{Introducción a la Lógica Proposicional (45 puntos)}
% ============================================================

\begin{enumerate}[leftmargin=*]
  \item \textbf{(8 pts)} Determina cuáles de los siguientes enunciados son proposiciones y asígnales un valor de verdad (V o F) cuando sea posible:
        \begin{enumerate}[label=\alph*)]
          \item «7 es un número primo.»
          \item «\textquestiondown Pásame el borrador?»
          \item «$x + 3 = 10$»
          \item «El Sol es un planeta.»
          \item «Todo número par es divisible por 2.»
          \item «\textquestiondown Cuántos años tienes?»
        \end{enumerate}

  \item \textbf{(8 pts)} Define las proposiciones atómicas necesarias y traduce al lenguaje simbólico:
        \begin{enumerate}[label=\alph*)]
          \item «Estudio y apruebo el examen.»
          \item «Si no llueve, entonces voy a la playa.»
          \item «Ni hace frío ni llueve.»
          \item «Llueve o hace sol.»
        \end{enumerate}

  \item \textbf{(9 pts)} Construye la \textbf{tabla de verdad completa} de la proposición:
        \[
          (p \vee q) \wedge \neg p
        \]

  \item \textbf{(10 pts)} En la isla de los caballeros (siempre dicen la verdad) y los pícaros (siempre mienten), te encuentras con dos habitantes $A$ y $B$. El habitante $A$ dice: «Ambos somos pícaros.» Determina qué tipo es cada uno. Justifica paso a paso.

  \item \textbf{(10 pts)} Tienes tres cajas etiquetadas «Manzanas», «Naranjas» y «Mixta», pero \textbf{todas las etiquetas están mal colocadas}. Solo puedes sacar \textbf{una fruta de una caja}. \textquestiondown De qué caja sacas la fruta y cómo determinas el contenido correcto de las tres cajas? Explica tu razonamiento paso a paso.
\end{enumerate}

\newpage

% ============================================================
\section{El Método de George Pólya (30 puntos)}
% ============================================================

\noindent Para cada problema, aplica explícitamente las cuatro fases del Método de Pólya (comprender, planificar, ejecutar y verificar).

\begin{enumerate}[leftmargin=*]
  \item \textbf{(7 pts)} En una granja hay gallinas y conejos. Se cuentan $25$ cabezas y $70$ patas. \textquestiondown Cuántas gallinas y cuántos conejos hay?

  \item \textbf{(7 pts)} Tienes billetes de $5$ y de $20$ bolívares. En total son $15$ billetes que suman $180$ bolívares. \textquestiondown Cuántos billetes de cada tipo tienes?

  \item \textbf{(8 pts)} Un tren sale de la ciudad $A$ hacia $B$ a $60$\,km/h. Dos horas después, otro tren sale de $B$ hacia $A$ a $90$\,km/h. Si las ciudades están a $420$\,km de distancia, \textquestiondown después de cuántas horas (desde que salió el segundo tren) se encuentran?

  \item \textbf{(8 pts)} En una prueba de $30$ preguntas, cada respuesta correcta vale $5$ puntos y cada incorrecta resta $2$ puntos (no hay preguntas sin responder). Un estudiante obtuvo $80$ puntos. \textquestiondown Cuántas respuestas correctas e incorrectas tuvo?
\end{enumerate}

\newpage

% ============================================================
\section{Del Lenguaje Natural al Lenguaje Algebraico (25 puntos)}
% ============================================================

\begin{enumerate}[leftmargin=*]
  \item \textbf{(5 pts)} Traduce a una expresión o ecuación algebraica (define tus variables):
        \begin{enumerate}[label=\alph*)]
          \item «El cuádruple de un número, disminuido en 7.»
          \item «La tercera parte de un número, más su cuadrado.»
          \item «La suma de dos números consecutivos es 51.»
        \end{enumerate}

  \item \textbf{(5 pts)} «La edad de Ana es el triple de la edad de su hijo. Juntas suman 60 años.» Plantea la ecuación y determina ambas edades.

  \item \textbf{(5 pts)} Un artículo cuesta $P$ bolívares y se le aplica un descuento del $20\%$.
        \begin{enumerate}[label=\alph*)]
          \item Escribe el precio final en función de $P$.
          \item Si el precio final es $640$ bolívares, \textquestiondown cuál era el precio original $P$?
        \end{enumerate}

  \item \textbf{(5 pts)} Un rectángulo tiene el largo $5$\,cm mayor que el ancho y su perímetro es $46$\,cm. Plantea la ecuación y halla las dimensiones.

  \item \textbf{(5 pts)} Escribe la identidad algebraica: «El cuadrado de la suma de dos números es igual a la suma de sus cuadrados más el doble de su producto.»
\end{enumerate}

\newpage

% ============================================================
\ifmwclaves
  \section{Clave de Respuestas}
  % ============================================================

  \subsection*{Sección 1: Introducción a la Lógica Proposicional}

  \begin{enumerate}[leftmargin=*, label=\textbf{\arabic*.}]
    \item a) Proposición, \textbf{V}. \quad b) No es proposición (es una orden/petición).
          \quad c) No es proposición (depende del valor de $x$). \quad d) Proposición, \textbf{F}.
          \quad e) Proposición, \textbf{V}. \quad f) No es proposición (es una pregunta).

    \item a) $p$: «Estudio», $q$: «Apruebo el examen» $\Rightarrow p \wedge q$. \quad
          b) $p$: «Llueve», $q$: «Voy a la playa» $\Rightarrow \neg p \rightarrow q$. \quad
          c) $p$: «Hace frío», $q$: «Llueve» $\Rightarrow \neg p \wedge \neg q$. \quad
          d) $p$: «Llueve», $q$: «Hace sol» $\Rightarrow p \vee q$.

    \item Tabla de verdad de $(p \vee q) \wedge \neg p$:
          \begin{center}
            \begin{tabular}{cc|c|c|c}
              \toprule
              $p$ & $q$ & $\neg p$ & $p \vee q$ & $(p \vee q) \wedge \neg p$ \\
              \midrule
              V   & V   & F        & V          & F                          \\
              V   & F   & F        & V          & F                          \\
              F   & V   & V        & V          & V                          \\
              F   & F   & V        & F          & F                          \\
              \bottomrule
            \end{tabular}
          \end{center}
          Se observa que es verdadera solo cuando $p$ es falsa y $q$ es verdadera.

    \item Supongamos que $A$ es caballero: entonces su afirmación sería verdadera, es decir, ambos serían pícaros, lo que incluiría al propio $A$: contradicción.
          Por lo tanto, $A$ es \textbf{pícaro} y su afirmación es falsa: no es cierto que ambos sean pícaros.
          Como $A$ ya es pícaro, necesariamente $B$ no lo es. Conclusión: $A$ es pícaro y $B$ es \textbf{caballero}.

    \item Se saca una fruta de la caja etiquetada \textbf{«Mixta»}. Como la etiqueta está mal, esa caja no contiene fruta mixta: contiene un solo tipo.
          \begin{itemize}[leftmargin=*]
            \item Si se saca una \textbf{manzana}, la caja «Mixta» es en realidad «Manzanas».
            \item Quedan las cajas etiquetadas «Manzanas» y «Naranjas», con contenidos Naranjas y Mixta (en algún orden). La caja «Naranjas» está mal etiquetada, así que no puede contener naranjas: es «Mixta». Por lo tanto, la caja «Manzanas» contiene «Naranjas».
            \item (Si la fruta sacada fuera una naranja, el razonamiento es simétrico.)
          \end{itemize}
  \end{enumerate}

  \newpage

  \subsection*{Sección 2: El Método de George Pólya}

  \begin{enumerate}[leftmargin=*, label=\textbf{\arabic*.}]
    \item \textbf{Fase 1:} $25$ animales, $70$ patas. Gallinas: 2 patas; conejos: 4 patas.
          \textbf{Fase 2:} suposición extrema: si todos fueran gallinas, habría $25 \times 2 = 50$ patas; faltan $70 - 50 = 20$. Cada conejo agrega $2$ patas.
          \textbf{Fase 3:} conejos $= 20 \div 2 = 10$; gallinas $= 25 - 10 = 15$.
          \textbf{Fase 4:} $15(2) + 10(4) = 30 + 40 = 70$ patas $\checkmark$ y $15 + 10 = 25$ cabezas $\checkmark$.
          \textbf{Respuesta:} 15 gallinas y 10 conejos.

    \item $x$: billetes de $5$; $y$: billetes de $20$. Sistema: $x + y = 15$, $5x + 20y = 180$.
          $x = 15 - y \Rightarrow 5(15 - y) + 20y = 180 \Rightarrow 75 + 15y = 180 \Rightarrow y = 7$, $x = 8$.
          \textbf{Verificación:} $8(5) + 7(20) = 40 + 140 = 180$ $\checkmark$. \quad 8 billetes de 5 y 7 de 20.

    \item \textbf{Comprender:} dos trenes que se acercan; el segundo sale $2$\,h después.
          \textbf{Plan:} al salir el segundo tren, el primero ya recorrió $60 \times 2 = 120$\,km; la distancia restante es $420 - 120 = 300$\,km. Se acercan a $60 + 90 = 150$\,km/h.
          \textbf{Ejecutar:} $t = \dfrac{300}{150} = 2$ horas después de que sale el segundo tren.
          \textbf{Verificar:} primer tren: $60 \times 4 = 240$\,km; segundo tren: $90 \times 2 = 180$\,km; total $240 + 180 = 420$\,km $\checkmark$.

    \item $c$: correctas; $i$: incorrectas. $c + i = 30$ y $5c - 2i = 80$.
          $i = 30 - c \Rightarrow 5c - 2(30 - c) = 80 \Rightarrow 5c - 60 + 2c = 80 \Rightarrow 7c = 140 \Rightarrow c = 20$, $i = 10$.
          \textbf{Verificación:} $5(20) - 2(10) = 100 - 20 = 80$ $\checkmark$. \quad 20 correctas y 10 incorrectas.
  \end{enumerate}

  \newpage

  \subsection*{Sección 3: Del Lenguaje Natural al Lenguaje Algebraico}

  \begin{enumerate}[leftmargin=*, label=\textbf{\arabic*.}]
    \item a) $4x - 7$ (con $x$ el número). \quad b) $\dfrac{x}{3} + x^2$. \quad c) $n + (n+1) = 51$.

    \item Sea $h$ la edad del hijo; Ana tiene $3h$. \quad $h + 3h = 60 \Rightarrow 4h = 60 \Rightarrow h = 15$.
          El hijo tiene $15$ años y Ana tiene $45$ años.

    \item a) Precio final $= P - 0{,}20P = 0{,}80P$. \quad b) $0{,}80P = 640 \Rightarrow P = \dfrac{640}{0{,}80} = 800$ bolívares.

    \item Sea $a$ el ancho; el largo es $a + 5$. Perímetro: $2\big(a + (a+5)\big) = 46 \Rightarrow 2a + 5 = 23 \Rightarrow a = 9$.
          \textbf{Dimensiones:} ancho $= 9$\,cm y largo $= 14$\,cm.

    \item $(a + b)^2 = a^2 + 2ab + b^2$.
  \end{enumerate}

\fi
\end{document}
